\documentclass[10pt,a4paper]{amsart}
\usepackage[margin=19mm]{geometry}
\usepackage{amsmath,amssymb,mathtools}
\usepackage[scr=rsfs,cal=euler]{mathalfa}
\usepackage{xfrac}
\usepackage[unicode,hidelinks,bookmarks=false]{hyperref}
\newcommand{\ie}{i.e.\ }
\newcommand{\cf}{cf.\ }
\newcommand{\resp}{respectively\ }
\newcommand{\Subsection}{\subsection}
\providecommand{\nobreakdash}{\nobreak\hbox{-}}
\setlength{\emergencystretch}{2em}
\multlinegap0pt
\usepackage{amssymb}
\usepackage{mathtools}
\usepackage[shortlabels]{enumitem}
\newtheorem{theorem}{Theorem}
\newcommand{\R}{\mathbb{R}}
\title{Radial Integration and Ball Volume\\ in Continuous Dimension}
\author{Andreu Ballus Santacana}
\date{Source: arXiv:2605.04351v1 (CC BY 4.0)}
\begin{document}
\maketitle

\noindent\emph{Source and licence.} Radial Integration and Ball Volume\\ in Continuous Dimension. Version arXiv:2605.04351v1 (CC BY 4.0), distributed by arXiv under the Creative Commons Attribution 4.0 International licence, http://creativecommons.org/licenses/by/4.0/, as recorded in the arXiv licence metadata for this exact version and shown on the abstract page https://arxiv.org/abs/2605.04351v1. Full licence notice: \texttt{licenses/arxiv-2605.04351-CC-BY-4.0.txt}.\par\smallskip

\noindent\textit{Editorial scope.} Counted unit: the article's theorem thm:homogeneous (source0.tex lines 563--569). The article's complete original proof of it is delivered with it (source0.tex lines 571--588, one \texttt{proof} environment of 18 lines). No mathematics is written, repaired or re-derived: the statement and the proof are the article's own lines, in the article's own notation, and no retained line is substituted or re-flowed.  Retained uncounted as the source-local context the proof needs: lines 554--556.  Nothing was omitted from any retained span, so the package carries no omission record.  No retained line contains a citation, so the wrapper carries no bibliography and no bibliography entry.  No retained line contains a figure environment, an image-inclusion command or drawing code, so no image is delivered and the package declares no asset dependency.  Editorial formatting only, in addition: an AMS \texttt{amsart} preamble that restates the article's own declarations and macros used by the retained lines, namely the environments \texttt{theorem} on separate counters, together with the standard packages those lines need.  The retained environments print in this document as Theorem 1 in order of appearance, whereas the article prints the same items with the numbering of its own full document.  The complete native source of the article as distributed in the arXiv e-print for this exact version is bundled unchanged as \texttt{sources/199851/source0.tex}.
\begin{equation}\label{eq:bui-randles-polar}
\int_{\R^d} f(\xi)\,d\xi \;=\; \int_0^\infty\!\int_{S_P} f(r^E\eta)\,d\sigma_P(\eta)\,r^{\mathrm{tr}(E)-1}\,dr
\end{equation}

\begin{theorem}[Homogeneous radial Gamma identity]\label{thm:homogeneous}
Under the hypotheses above, the unit sublevel set $B_P := \{\xi\in\R^d : P(\xi)<1\}$ has Lebesgue measure
\begin{equation}\label{eq:homogeneous-gamma}
m(B_P) \;=\; \frac{1}{\Gamma(\mu_P+1)}\int_{\R^d} e^{-P(\xi)}\,d\xi.
\end{equation}
Moreover, $\sigma_P(S_P)\,\Gamma(\mu_P) = \int_{\R^d} e^{-P(\xi)}\,d\xi$, so~\eqref{eq:homogeneous-gamma} is equivalent to $m(B_P) = \sigma_P(S_P)/\mu_P$.
\end{theorem}

\begin{proof}
By~\eqref{eq:bui-randles-polar} applied to $f(\xi) := \Phi(P(\xi))$ for any Borel $\Phi\colon[0,\infty)\to[0,\infty)$, and using $P(r^E\eta) = r$ for $\eta\in S_P$,
\[
\int_{\R^d}\Phi(P(\xi))\,d\xi \;=\; \int_0^\infty\!\int_{S_P}\Phi(r)\,d\sigma_P(\eta)\,r^{\mu_P-1}\,dr \;=\; \sigma_P(S_P)\int_0^\infty \Phi(r)\,r^{\mu_P-1}\,dr.
\]
Take $\Phi(r) = e^{-r}$:
\begin{equation}\label{eq:homog-gauss}
\int_{\R^d} e^{-P(\xi)}\,d\xi \;=\; \sigma_P(S_P)\int_0^\infty e^{-r}\,r^{\mu_P-1}\,dr \;=\; \sigma_P(S_P)\,\Gamma(\mu_P).
\end{equation}
Take $\Phi(r) = \mathbf 1_{(0,1)}(r)$:
\begin{equation}\label{eq:homog-ball}
m(B_P) \;=\; \int_{\R^d}\mathbf 1_{(0,1)}(P(\xi))\,d\xi \;=\; \sigma_P(S_P)\int_0^1 r^{\mu_P-1}\,dr \;=\; \frac{\sigma_P(S_P)}{\mu_P}.
\end{equation}
Eliminating $\sigma_P(S_P)$ between~\eqref{eq:homog-gauss} and~\eqref{eq:homog-ball}, and using $\mu_P\,\Gamma(\mu_P)=\Gamma(\mu_P+1)$,
\[
m(B_P) \;=\; \frac{1}{\mu_P\,\Gamma(\mu_P)}\int_{\R^d}e^{-P(\xi)}\,d\xi \;=\; \frac{1}{\Gamma(\mu_P+1)}\int_{\R^d}e^{-P(\xi)}\,d\xi. \qedhere
\]
\end{proof}

\end{document}
